% 解説プリントのコメント（saiten が書く。./tools/kaisetsu.ps1 がこのファイルからプリントを作る）
%% matome
\begin{daiji}
  \item 複素数を成分にもつベクトルの大きさの2乗は、各成分の絶対値の2乗の和である。$\abs{1+\ii}^{2}=(1+\ii)(1-\ii)=2$ で、4 ではない（問題1 (2)）。
  \item 「エルミート行列なら固有値は実数」の逆は言えない。エルミートでない行列の固有値は「実数とは限らない」のであって「実数でない」のではない（問題1 (4)）。
\end{daiji}
\begin{maemo}
  \item 問題1 (4) は同じ問題で同じ誤り：2026-10-28（当時の第4回 問題1 (4)）にも「エルミートでないため実数にならない」と書いた。今回も「エルミートでないので固有値は実数でない」と書いている。答えの文は「$Q$ はエルミートでないので固有値が実数であるとは保証されず、実際 $1\pm\ii$ である」と、「保証されない」の形で書く。
  \item 複素数の大きさの2乗の扱い：量子力学1-3 問題2（2026-10-18・10-21）で $\abs{A}^{2}$ を $A^{2}$、量子力学1-13 問題2・3（2026-11-21）で $\abs{r}^{2}$ を $r^{2}$ と書いた。今回は $\abs{1-\ii}^{2}$ を 4 とした。規格化したら、最後に $\vect{u}^{\dagger}\vect{u}=1$ を計算して確かめる。
\end{maemo}
%% mondai 1
\begin{kotaeran}
  \machigai{2}{△}{計算} 固有値 $\lambda=0,\,3$ と、成分の比 $(1+\ii,\,-1)^{\mathrm{T}}$・$(1,\,1-\ii)^{\mathrm{T}}$ は正しい（どちらも両方の行の式を満たす）。規格化の係数を両方とも $1/\sqrt{5}$ とした。大きさの2乗は $\abs{1+\ii}^{2}+\abs{-1}^{2}=2+1=3$、$1+\abs{1-\ii}^{2}=1+2=3$ なので、正しくは $1/\sqrt{3}$。$\abs{1\pm\ii}^{2}$ を 4 としたことになる。また、第1行の式から決めた比が第2行の式も満たすことの確認が書かれていない（解答の論点）。(3) の係数 $1/5$ はこの誤りを引き継いだもので、内積が 0 になる結論には影響しない。
  \machigai{4}{○}{記述} 固有値 $1\pm\ii$ は正しい。理由を「$Q$ はエルミートでないので固有値は実数でない」と書いた。定理は「エルミートなら実数」であり、エルミートでない行列でも実数の固有値をもつものはある（例えば $\begin{pmatrix}1&1\\0&2\end{pmatrix}$ の固有値は 1 と 2）。正しくは「$Q$ はエルミートでなく定理の前提を満たさないので、固有値が実数でなくても矛盾しない」。
\end{kotaeran}
\begin{rikai}
  \item 複素ベクトル $(a,b)^{\mathrm{T}}$ の大きさの2乗は $\abs{a}^{2}+\abs{b}^{2}=a^{*}a+b^{*}b$。$\abs{x+\ii y}^{2}=x^{2}+y^{2}$ なので $\abs{1\pm\ii}^{2}=2$。規格化の後に $\vect{u}^{\dagger}\vect{u}=\frac{1}{3}(2+1)=1$ を計算すれば、係数の誤りはすぐ見つかる。
  \item 「$A$ ならば $B$」から言えるのは「$B$ でないならば $A$ でない」（対偶）であって、「$A$ でないならば $B$ でない」（裏）ではない。エルミートでない行列については、固有値が実数であることが「保証されない」とだけ言える。
\end{rikai}
